%% gyoudori-en.tex: manual of the gyoudori package (English)
%% Build: platex gyoudori-en && platex gyoudori-en && dvipdfmx gyoudori-en
\documentclass[paper=a4,dvipdfmx]{jlreq}
\usepackage{amsmath}
\usepackage{url}
\usepackage{tikz}
\usetikzlibrary{decorations.pathreplacing}
\usepackage[displaymath]{gyoudori}
%% gyoudori-fig.tex: マニュアルの図（日本語版・英語版で共通）
%% 寸法は模式的：字の枠10，行間5，行送り15（単位pt）．ベースラインは字の枠の上端から8.8

%% 行取りの領域
%% #1：本文，#2：字の枠，#3：行間，#4：行送り，#5：3行取りの領域，
%% #6：見出し，#7：領域の高さの式
\newcommand\GyFigArea[7]{%
  \begin{tikzpicture}[x=1.5pt,y=1.8pt,font=\footnotesize]
    % 本文の行（1，2行目と6，7行目）
    \foreach \i in {0,1,5,6} {
      \fill[black!15] (0,-\i*15) rectangle (150,-\i*15-10);
    }
    \node[anchor=west] at (152,-5) {#1};
    \node[anchor=west] at (152,-80) {#1};
    % 行取りの領域の中の行の位置（3から5行目）
    \foreach \i in {2,3,4} {
      \draw[densely dotted] (0,-\i*15) rectangle (150,-\i*15-10);
    }
    % 領域
    \draw[thick] (0,-30) rectangle (150,-70);
    % 見出し（領域の中央）
    \fill[black!45] (20,-42) rectangle (120,-58);
    \node[white,font=\small\bfseries] at (70,-50) {#6};
    % 右：領域の高さ
    \draw[decorate,decoration={brace,amplitude=4pt}] (156,-30) -- (156,-70)
      node[midway,anchor=west,xshift=6pt,align=left] {#5\\#7};
    % 左：字の枠，行間，行送り
    \draw[|-|] (-4,0) -- (-4,-10)
      node[midway,anchor=east,xshift=-2pt] (fr) {#2};
    \draw[|-|] (-4,-10) -- (-4,-15) node[midway,anchor=east,xshift=-2pt] {#3};
    % 行送り：ベースラインどうしの間隔（このパッケージの基準点）．
    % 寸法線は「字の枠」のラベルの左に置く
    \coordinate (P) at ([xshift=-7pt]fr.west);
    \foreach \y in {-8.8,-23.8} {
      \draw[densely dashed,very thin] ([xshift=-5pt]P |- 0,\y) -- (150,\y);
    }
    \draw[|-|] (P |- 0,-8.8) -- (P |- 0,-23.8) node[midway,anchor=east,xshift=-2pt] {#4};
  \end{tikzpicture}}

%% 領域の中の位置（t/c/b）
%% #1：本文，#2から#4：t，c，bの説明
\newcommand\GyFigPos[4]{%
  \begin{tikzpicture}[x=1.3pt,y=1.3pt,font=\footnotesize]
    \foreach \col/\lab/\ya/\yb in {0/#2/-30/-46, 90/#3/-42/-58, 180/#4/-54/-70} {
      \begin{scope}[shift={(\col,0)}]
        \fill[black!15] (0,0) rectangle (70,-10);
        \fill[black!15] (0,-15) rectangle (70,-25);
        \foreach \i in {2,3,4} {
          \draw[densely dotted] (0,-\i*15) rectangle (70,-\i*15-10);
        }
        \draw[thick] (0,-30) rectangle (70,-70);
        \fill[black!45] (8,\ya) rectangle (62,\yb);
        \fill[black!15] (0,-75) rectangle (70,-85);
        \node at (35,-95) {\lab};
      \end{scope}
    }
  \end{tikzpicture}}

%% 3行取りの領域がページの途中と先頭に来たとき
%% 4つとも版面の天から7行分を描き，左に行番号をつける
%% #1：見出し，#2：ページの途中（中央），#3から#5：ページの先頭（上寄せ，中央，下寄せ），
%% #6：版面の天
\newcommand\GyFigTop[6]{%
  \begin{tikzpicture}[x=1.12pt,y=1.2pt,font=\footnotesize]
    % ##1：段の左端，##2：領域の1行目（0始まり），##3：見出しの行，##4：下の説明
    \def\GyPanel##1##2##3##4{%
      \begin{scope}[shift={(##1,0)}]
        \draw[thick] (-5,4) -- (65,4);
        \draw (-5,4) -- (-5,-104); \draw (65,4) -- (65,-104);
        \node[anchor=south west,inner xsep=0pt] at (-5,5) {#6};
        \foreach \i in {0,...,6} {
          \node[font=\tiny] at (-10,-\i*15-5) {\pgfmathparse{int(\i+1)}\pgfmathresult};
          \pgfmathparse{(\i<##2)||(\i>##2+2)}
          \ifnum\pgfmathresult=1
            \fill[black!15] (0,-\i*15) rectangle (60,-\i*15-10);
          \else
            \draw[densely dotted] (0,-\i*15) rectangle (60,-\i*15-10);
          \fi
        }
        \draw[thick] (0,-##2*15) rectangle (60,-##2*15-40);
        \fill[black!45] (6,-##3*15) rectangle (54,-##3*15-10);
        \node[white,font=\scriptsize\bfseries] at (30,-##3*15-5) {#1};
        \node[align=center] at (30,-118) {##4};
      \end{scope}}
    \GyPanel{0}{2}{3}{#2}
    \GyPanel{95}{0}{0}{#3}
    \GyPanel{190}{0}{1}{#4}
    \GyPanel{285}{0}{2}{#5}
  \end{tikzpicture}}

%% 実例の背景に，行の位置（字の枠）をグレーの帯で描く
%% 段落の1行目の先頭に置く．#1：帯の数．寸法はgyoudoriと同じ（\normalsizeの
%% 「あ」の高さ・深さと\normalbaselineskip）
\makeatletter
\newcommand\GyBands[1]{%
  \leavevmode
  \@gyoudori@metrics
  \edef\GyHt{\strip@pt\@gyoudori@ht}%
  \edef\GyDp{\strip@pt\@gyoudori@dp}%
  \edef\GyBs{\strip@pt\@gyoudori@bs}%
  \edef\GyWd{\strip@pt\linewidth}%
  \rlap{\begin{tikzpicture}[overlay,x=1pt,y=1pt]
    \foreach \i in {1,...,#1} {
      \fill[black!12] (0,\GyHt-\i*\GyBs+\GyBs) rectangle
        (\GyWd,-\GyDp-\i*\GyBs+\GyBs);
    }
  \end{tikzpicture}}}
\makeatother


\newcommand\pkg[1]{\textsf{#1}}
\newcommand\opt[1]{\texttt{#1}}
\newcommand\env[1]{\texttt{#1}}
\newcommand\cs[1]{\texttt{\char`\\#1}}
\newcommand\meta[1]{$\langle$\textit{#1}$\rangle$}
\newenvironment{demo}
  {\par\medskip\noindent
   \begin{minipage}{\linewidth}\hrule\smallskip}
  {\par\smallskip\hrule\end{minipage}\par\medskip}
\renewcommand\figurename{Figure\nobreakspace}
\renewcommand\abstractname{Abstract}
\renewcommand\refname{References}

\title{The \pkg{gyoudori} package\\[.5ex]\large Gyou-dori (line-based placement)
  for Japanese typesetting}
\author{Yuwsuke Kieda}
\date{2026-10-16\quad v20261016.0}

\begin{document}
\maketitle

\begin{abstract}
The \pkg{gyoudori} package places headings, display math and floats by
\emph{gyou-dori}, so that the lines of the main text before and after them stay
on the line positions of the kihon-hanmen (the line grid). It works with pLaTeX,
upLaTeX and LuaLaTeX (LuaTeX-ja). A Japanese manual (\texttt{gyoudori-ja.pdf})
is also available.
\end{abstract}

\section{What is gyou-dori?}

\emph{Gyou-dori} (行取り, also written 行ドリ) is one of the arbitrary layouts: the
space in the block direction occupied by an element whose size differs from the
main text, such as a heading, is determined as a number of lines of the
kihon-hanmen (the basic page format), so that the lines of the main text after
it are not shifted off the line positions defined by the kihon-hanmen (the line
grid).

Gyou-dori is used in typesetting Japanese books and the like when it is
specified that the lines of the main text on each page be placed on the line
grid. It is called arbitrary because aligning lines to the grid is not in itself
a principle of readability. A grid layout is often specified regardless of, or
even ignoring, the benefit of readers, and is not seldom a choice that abandons
the consideration of page design. Page design should rather pursue the benefit
of readers without being tightly bound to the line grid. This package is a tool
for cases where one has to typeset under such a specification.

\subsection{The gyou-dori area}

As defined in JLREQ (W3C \textit{Requirements for Japanese Text Layout})
\S4.1.6, the block-direction size of the area of an
$n$-line gyou-dori is
\[
  \text{character frame}\times n + \text{line gap}\times(n-1)
  = \text{line pitch}\times(n-1) + \text{character frame},
\]
that is, the range from the top of the character frame of the first line to the
bottom of the character frame of the $n$-th line (Figure~\ref{fig:area}).
Visually, a line gap is added before and after the area in the middle of a page
(or column), and only after it at the top of a page (or column).
The line pitch in the figure is shown as the distance between baselines (dashed
lines). In this manual, the reference point for measuring the line pitch is
tentatively taken to be the baseline (Section~\ref{sec:pitch}).

\begin{figure}[htbp]
  \centering
  \GyFigArea{text}{char.\ frame}{line gap}{line pitch}{3-line area}{heading}%
    {$3\times\text{frame}+2\times\text{gap}$}
  \caption{A 3-line gyou-dori area with a heading at its center}
  \label{fig:area}
\end{figure}

By default this package places the content at the center of the area. The
position argument (\opt{t}, \opt{c}, \opt{b}) changes it to top or bottom
alignment (Figure~\ref{fig:pos}). Where in the area the content should be placed
depends on the specification of each typesetting job; JLREQ, too, lists methods
other than centering (Section~\ref{sec:pos}).

\begin{figure}[htbp]
  \centering
  \GyFigPos{text}{\opt{t}: top}{\opt{c}: center (default)}{\opt{b}: bottom}
  \caption{Positions in the area}
  \label{fig:pos}
\end{figure}

\subsection{Dimensions used by this package}\label{sec:dims}

The line pitch is \cs{normalbaselineskip}, and the height of the character
frame is the height plus depth of the hiragana ``あ'' (the JFM metrics), both taken
at \cs{normalsize}. Even inside \cs{small} etc., the grid of the main text is
used (except with \opt{local} for the \env{gyoudori} environment and
\cs{gyoudoriarea}; Section~\ref{sec:env}).

The height of the character frame depends on the JFM. With the standard JFMs of
upLaTeX and LuaTeX-ja (including \pkg{jlreq}) it equals the full-width size (a
square frame), which agrees with JLREQ. With the standard JFM of pLaTeX
(\texttt{jis}) it is 0.952 times the full-width size. Since this package
follows the frame of the JFM, an area in pLaTeX is shorter than the JLREQ
definition by 0.048 times the full-width size.

\subsection{Reference point of line pitch and \cs{baselineskip}}\label{sec:pitch}

JLREQ (note in \S2.4.2) defines the line feed (line pitch) as ``the distance
between two adjacent lines measured from their reference points'', but does not
fix the reference point itself. It says that the reference point differs from
implementation to implementation and that the vertical center of the character
frame is commonly used in horizontal writing, without giving any evidence for
this. Moreover, \LaTeX{} (\TeX{}) stacks lines by their baselines, which is far
from the idea of using the vertical center as the reference point.

This manual therefore tentatively takes the baseline as the reference point of
the line pitch, because both the default behavior of \LaTeX{} and the
implementation of this package are based on baselines. That is, the line pitch
is the distance between baselines.

\cs{baselineskip} of \LaTeX{} is the distance between the baselines of adjacent
lines. \TeX{} inserts glue between lines, computed from the depth of the previous
line and the height of the next, so that the baselines are \cs{baselineskip}
apart; however, if the space between the lines would be less than
\cs{lineskiplimit}, \cs{lineskip} is inserted instead and the distance becomes
larger than \cs{baselineskip}.

This package uses \cs{normalbaselineskip} at \cs{normalsize} as the line pitch
$b$ of the main text (a change by \cs{linespread} is reflected in it), and
regards the positions where baselines are $b$ apart as the line positions of the
main text (the line grid). The character frame of each line extends from the
baseline upward by the height $h$ of ``あ'' and downward by its depth $d$. Hence
the area of an $n$-line gyou-dori extends from $h$ above the baseline of the first
line to $d$ below the baseline of the $n$-th line, and its size is $(n-1)b+h+d$.
The area is placed so that the baseline of its first line is at distance $b$ from
the baseline of the preceding line, and the baseline of the line after the area
is at distance $b$ from the baseline of the $n$-th line
(Section~\ref{sec:place}). Display math, too, is placed so that the midpoint
between the baselines of its first and last rows coincides with the midpoint
between the baselines of the first and last lines of the area
(Section~\ref{sec:math}).

Between lines consisting of characters of the same size, the value of the line
pitch does not depend on the choice of the reference point; neither does the
JLREQ relation (line feed $=$ character size $+$ line gap) nor the formula for
the size of the area. The reference point matters only for lines with
characters of different sizes, such as a large heading: then the distance
between baselines differs from that between vertical centers. This package puts
such headings into gyou-dori areas and keeps the surrounding lines of the main
text on the line grid by their baselines.

The line gap of this package (line pitch $-$ height of the character frame)
equals the line gap of JLREQ (line feed $-$ character size) for JFMs with a
square frame. With the standard JFM of pLaTeX, whose frame height is 0.952 times
the full-width size, it is larger than the JLREQ line gap by 0.048 times the
full-width size.

\section{Requirements and installation}

\begin{itemize}
\item pLaTeX, upLaTeX, or LuaLaTeX with LuaTeX-ja (\pkg{ltjsarticle},
  \pkg{jlreq}, etc.). Horizontal writing only.
\item \LaTeX{} 2020-10-01 or later (the package uses \cs{AddToHook}).
\item The \pkg{kvoptions} package.
\end{itemize}

Put \texttt{gyoudori.sty} into your TEXMF tree (e.g.\ \texttt{tex/latex/gyoudori/}).

\section{Usage}

\subsection{The \env{gyoudori} environment}\label{sec:env}

\begin{verbatim}
\begin{gyoudori}[<options>]{<lines>}
<content>
\end{gyoudori}
\end{verbatim}
places \meta{content} in a \meta{lines}-line gyou-dori area. \meta{options} is a
comma-separated list of a position and \opt{local}. The position is \opt{t}
(top), \opt{c} (center, default) or \opt{b} (bottom). The content is set in
vertical mode, so it may contain several lines or size changes.

With \opt{local}, the gyou-dori is based not on the line pitch of
\cs{normalsize} but on the height and depth of ``あ'' in the current font size
and on \cs{baselineskip}. Use it to keep the surrounding lines on their grid
inside \cs{small} etc.\ (e.g.\ \verb|\begin{gyoudori}[b,local]{2}|). Without
\opt{local}, the line pitch of \cs{normalsize} is used regardless of the font
size.

\begin{demo}
A line of text. This is the last line of the preceding paragraph.
\begin{gyoudori}{3}
\large A heading centered in 3 lines
\end{gyoudori}

The paragraph after the heading stays on the line grid.
\begin{gyoudori}[b]{2}
A heading at the bottom of 2 lines
\end{gyoudori}
Text that follows without a blank line is not indented.
\end{demo}

The next example uses \opt{local} inside a part set in \cs{small}. The gray
bands show the line positions of \cs{small}.

\begin{demo}
\small\GyBands[local]{5}A line of text set in \cs{small}.
\begin{gyoudori}[local]{3}
\normalsize\bfseries A heading centered in 3 lines
\end{gyoudori}
The following line also stays on the line positions of \cs{small}.
\end{demo}

\begin{itemize}
\item If no blank line follows \verb|\end{gyoudori}|, the following text is
  neither indented nor preceded by \cs{parskip}, as after a display math.
  A blank line starts a new paragraph, indented and preceded by \cs{parskip}.
\item The area is never split. If it does not fit at the bottom of a page
  (column), the whole heading moves to the top of the next page (column)
  (JLREQ \S4.1.7~a, c).
\item No page (column) break is allowed right after the heading. If not even
  one line of text fits after the heading, the whole heading moves to the next
  page (column) (JLREQ \S4.1.7~b). In both cases the bottom of the previous page
  (column) is left blank.
\end{itemize}

\subsection{The \cs{gyoudoriarea} command}

\begin{verbatim}
\gyoudoriarea[<options>]{<lines>}{<content>}
\end{verbatim}
works like the \env{gyoudori} environment, with the same \meta{options}. Being a command, it does not control
the indentation of the following text.

\subsection{Placement methods of JLREQ}\label{sec:pos}

JLREQ \S4.1.6 lists four methods (a--d) of placing a heading by gyou-dori.
They are written with this package as follows.
\begin{description}
\item[a. Center of a multi-line area] \verb|\begin{gyoudori}{3}|.
\item[b. Centered, with blank lines before or after] Combine with a space of
  whole lines. For a one-line space before the heading, put
  \verb|\vspace{\baselineskip}| before the environment.
\item[c. A specified position in the area] Top (\opt{t}) and bottom (\opt{b})
  are provided. For other positions, put space into the content.
\item[d. A one-line heading on a line position with a blank line before] For a
  heading in the text size, a 2-line area with bottom alignment (position
  \opt{b}, 2 lines) gives the same result.
\end{description}
\paragraph{When a heading area comes to the top of a page}
The area of a heading, too, is placed at the top of the kihon-hanmen at the top
of a page (or column), and the position in the area follows the position option
(Section~\ref{sec:edge}). For example, when the ``low-level heading with a
one-line space before it'' of method d (a 2-line area with bottom alignment)
comes to the top of a page, the first line is left blank and the heading is on
the second line. This is the policy ``always keep the one-line space'' of JLREQ
\S4.1.8.

\subsection{Gyou-dori areas at the top and bottom of a page}\label{sec:edge}

Not only headings: an area of the \env{gyoudori} environment, of
\cs{gyoudoriarea}, or of display math may happen to come to the top or bottom of
a page (or column) as the text flows. It then behaves as follows. Floats are not
covered here, since \LaTeX{} decides where they go (Section~\ref{sec:float}).

\paragraph{When a gyou-dori area comes to the top of a page}
The area itself is placed at the top of the kihon-hanmen (\emph{tentsuki}), with
no line gap before it. The position in the area follows the position option
(\opt{t}, \opt{c}, \opt{b}) as in the middle of a page
(Figure~\ref{fig:top}), and the line after an $n$-line area starts on line
$n+1$. The same holds for display math: the formula is placed in the area at the
same position as in the middle of a page.

\begin{figure}[htbp]
  \centering
  \GyFigTop{heading}{middle of a page\\(\opt{c})}{top of a page\\(\opt{t})}%
    {top of a page\\(\opt{c})}{top of a page\\(\opt{b})}{top of the kihon-hanmen}
  \caption{A 3-line area in the middle and at the top of a page (numbers on the
    left are line numbers)}
  \label{fig:top}
\end{figure}

\paragraph{When a gyou-dori area comes to the bottom of a page}
The area is never split. If it fits by the last line of the kihon-hanmen, it
stays on the page; otherwise the whole area moves to the top of the next page
(column), and the bottom of the previous page (column) is left blank for the
lines that did not fit.
\begin{itemize}
\item The \env{gyoudori} environment and \cs{gyoudoriarea} are meant for
  headings and allow no page (column) break right after them. If the area ends
  on the last line and not even one line of text fits after it, the whole area
  moves to the next page (column) (Section~\ref{sec:env}).
\item A page break is allowed after display math. If the area ends on the last
  line, it stays on the page and the following text starts on the first line of
  the next page (column). A multi-line display is a single area, so it is not
  broken across pages even with \cs{allowdisplaybreaks}.
\end{itemize}

\section{Package options}

Options are given as key-value pairs. For a Boolean option, the key alone means
\opt{true}.
\begin{verbatim}
\usepackage[displaymath,float,textfloatsep=2\baselineskip]{gyoudori}
\end{verbatim}

\subsection{\opt{displaymath=true|false}}

Display math is placed by gyou-dori automatically. The default is \opt{false}.
The environments \env{eqnarray}, \env{align} and \env{gather} and their starred
forms are handled (\env{align} and \env{gather} need \pkg{amsmath}). Each display
is placed so that the midpoint between the baselines of its first and last rows
coincides with the midpoint between the baselines of the first and last lines of
the area; a one-row display in an area of an odd number of lines thus sits on
the baseline of the middle line. The number of lines is the smallest (at least
two) with which the display fits at this position; the display may extend into
the outer line gaps by up to half of them, so that at least half of the line gap
remains to the character frames of the adjacent lines.
\cs{abovedisplayskip} and the like are not used. Equation numbers, \cs{label}
and \cs{eqref} work as they do without this option.

\begin{demo}
\GyBands{8}A line of text. The following display is placed by gyou-dori.
\begin{align}
  a &= b + c \\
  d &= e
\end{align}
The text after the display is on the grid.
\begin{gather*}
  x = y
\end{gather*}
A line of text.
\end{demo}

The gray bands show the line positions (the character frame of each line).
In this example, the \env{align} is centered in a 3-line area and the
\env{gather*} in a 2-line area, and the text before and after them stays on the
bands. Since \env{align} and \env{gather} of \pkg{amsmath} insert \cs{strut} in
every row, the number of lines of a two-row \env{align} depends on the line
pitch of the class (three lines with \pkg{jlreq} used for this manual, two lines
with \pkg{jsarticle}).

Larger displays with fractions, integrals and sums are also placed in the smallest area that holds them.

\begin{demo}
\GyBands{13}A line of text. An integral and a sum:
\begin{gather}
  \int_0^1 f(x)\,dx = \sum_{k=1}^{n} \frac{a_k}{k!}
\end{gather}
Text after the display. Two rows:
\begin{align}
  S_n &= \sum_{k=1}^{n} k = \frac{n(n+1)}{2} \\
  I &= \int_{-\infty}^{\infty} e^{-x^2}\,dx = \sqrt{\pi}
\end{align}
More text. A case analysis:
\begin{gather}
  f(x) =
  \begin{cases}
    \frac{x^2}{2} & (x \ge 0) \\
    -x & (x < 0)
  \end{cases}
\end{gather}
The last line of text.
\end{demo}

In this example, the displays take two, four and three lines, and the text after each of them starts on a band.

\subsection{\opt{float=true|false}}\label{sec:float}

Floats are placed by gyou-dori. The default is \opt{false}.
\begin{itemize}
\item A float is put into the smallest gyou-dori area (at least one line) that
  holds it.
\item When its place is decided, the float is aligned in the area: at the top
  of a page, to the top, with the top of the area on the top of the character
  frame of the first line; at the bottom, to the bottom, with the bottom of the
  area on the bottom of the character frame of a line; and here (\opt{h}), to the
  center, aligned with the lines before and after it.
\item The five float separations (\cs{floatsep} etc.; see below) are rounded up
  to the smallest value of the form ``line gap $+$ a multiple of the line
  pitch'' that is not less than the given value, without stretch or shrink.
  \cs{intextsep}, which is used both above and below, takes half of the line gap.
\end{itemize}
Thus the text lines around floats stay on the grid. \LaTeX{} changes a lone
\opt{[h]} into \opt{[ht]}; a float that goes to the top instead is realigned to
the top. No column break is allowed right before a float placed here.

\subsection{\opt{floatsep}, \opt{textfloatsep}, \opt{intextsep},
  \opt{dblfloatsep}, \opt{dbltextfloatsep}}

Set the length of the same name. If not given, the value of the document class
is used. With \opt{float=true}, the value is rounded as described above.

\section{Use with other packages}

The following applies with \opt{float=true}.
\begin{description}
\item[\pkg{nidanfloat}] A double-column float (\env{figure*}) is aligned to the
  top of the page, or to the bottom of the character frame of the last line
  that fits in the page. A single-column bottom float in two-column mode is
  aligned to the last line of the column. The loading order does not matter.
\item[\pkg{stfloats}] Double-column floats at the top and bottom are aligned
  likewise.
\item[\pkg{flushend}] The glue in the balanced columns of the last page is not
  stretched or shrunk, so the lines stay on the grid (the \opt{spread} and
  \opt{shrink} options have no effect). The gap where the two columns are joined
  is rounded up to a multiple of the line pitch.
\item[\pkg{evenend}] (LuaLaTeX) Floats on the balanced last page are aligned to
  the top or bottom of the column.
\end{description}

\section{Relation to JLREQ}

The behavior of this package was checked against the description of gyou-dori
in JLREQ.

\begin{description}
\item[Reference point of the line pitch (note in \S2.4.2)] JLREQ does not fix it.
  This package takes the baseline as the reference point
  (Section~\ref{sec:pitch}).
\item[Size of the area (\S2.5.2, \S4.1.6)] Agrees. As the frame of the JFM is
  used, an area in pLaTeX with the standard JFM is shorter than the definition
  with a square frame by 0.048 times the full-width size (Section~\ref{sec:dims}).
\item[Placement in the area (\S4.1.6 a--d)] a is the default center and c is
  covered by top and bottom alignment; b and d are realized by combination
  (Section~\ref{sec:pos}).
\item[Heading at the bottom of a page or column (\S4.1.7)] Agrees. If the area
  does not fit, or if no text line fits after the heading, the whole heading
  moves to the next page (column). The exception for even pages in vertical
  writing does not apply, since only horizontal writing is supported.
\item[Heading at the top of a page (\S4.1.6, \S4.1.8)] Agrees. The area is
  placed at the top of the kihon-hanmen and only the line gap after it is added
  (\S4.1.6). For a low-level heading with a one-line space before it, this is the
  policy ``always keep the one-line space'' among the three in \S4.1.8
  (Section~\ref{sec:pos}).
\item[Blank space at the end of a column (\S4.1.7 c)] JLREQ prefers some
  processing to avoid it but allows it if unavoidable. This package leaves it
  blank.
\item[Display math (note in \S4.1.6)] For a display formula larger than the
  text size, JLREQ mentions two methods: specifying the area by gyou-dori, or
  keeping a minimum line gap before and after and adjusting the remainder within
  the page or column. \opt{displaymath} follows the former. The minimum of two
  lines is a design choice of this package.
\item[Illustrations (\S2.5.2, \S4.3)] JLREQ says that the block-direction space
  of illustrations should also be aligned to the line positions of the
  kihon-hanmen; \opt{float} follows this. Avoiding a single text line before or
  after an illustration (\S4.3) is left to the float placement of \LaTeX.
\end{description}

\section{How it works}

\subsection{Notation}

$b$: the line pitch (\cs{normalbaselineskip} at \cs{normalsize});
$h$, $d$: the height and depth of ``あ'' at \cs{normalsize} (the character
frame); $g=b-h-d$: the line gap; $A(n)=(n-1)b+h+d$: the height of an $n$-line
area.

\subsection{Placing a gyou-dori area}\label{sec:place}

The content is set in a \cs{vbox} of height $A(n)$ with \cs{vss} after, around
or before it (for \opt{t}, \opt{c} and \opt{b}). The height of this box is changed
to $h$ and its depth to $A(n)-h$, and the box is put after the paragraph is
ended. \TeX{} inserts \cs{baselineskip} glue so that the baseline of the box (the
baseline of the first line of the area) is at distance $b$ from the previous
baseline. After the box, \cs{prevdepth} is set to $d$, so the next line is at
distance $b$ from the baseline of the last line of the area, i.e.\ on line $n+1$.
At the top of a page (column), \cs{topskip} puts the top of the area on the top
of the character frame of the first line. This is the method of Noriyuki Abe
(\url{https://abenori.blogspot.com/2017/12/gyodori.html}).
With \opt{local}, the height and depth of ``あ'' in the current font size and
\cs{baselineskip} are used for $h$, $d$ and $b$.

The \env{gyoudori} environment and \cs{gyoudoriarea} add \cs{nobreak} after the
area to forbid a page break right after a heading. After the \env{gyoudori}
environment and the display math environments, the same mechanism as
\cs{@doendpe} of \LaTeX{} suppresses the indentation and sets \cs{parskip} to
zero for the next paragraph unless a blank line follows.

\subsection{Display math}\label{sec:math}

\begin{enumerate}
\item With the environment hooks (\texttt{env/<name>/before} and \texttt{after}),
  a \cs{vbox} is opened outside the environment, and the display is set in it
  with \cs{abovedisplayskip} etc.\ set to zero.
\item A \cs{penalty}~$-10000$ is put at the top of this box. When a copy of the
  finished box is split by \cs{vsplit} with \cs{splittopskip} set to a large value
  $S$, the rest begins with glue of ``$S-{}$the height of the first row''. Thus
  ``the height of the rest $-S$'' is the distance $\Delta$ between the baselines of
  the first and last rows.
\item Let $H$ and $D$ be the height and depth of the box. The baseline of the
  last row is at $H$ from the top of the box, and that of the first row at
  $H-\Delta$. With the distance from the top of the area to the top of the box
  \[ y = h + \frac{(n-1)b}{2} - H + \frac{\Delta}{2}, \]
  the midpoint of the baselines of the display coincides with that of the area.
\item Starting from $n=2$, the smallest $n$ with $-g/2\le y$ and
  $y+H+D\le A(n)+g/2$ is used (the box may extend into the outer line gaps by up
  to half of them).
\item The box is lowered by $y$ in the area, which is placed in the same way
  as in Section~\ref{sec:place}.
\end{enumerate}

\subsection{Floats}

\begin{description}
\item[Size] After \cs{@endfloatbox}, the float box is rebuilt as a \cs{vbox} of
  height $A(n)$ for the smallest $n\ge1$ with height $+$ depth $\le A(n)$ (the
  content at the top).
\item[Realignment] When the placement is decided, the content is placed at
  position $p$ (\opt{t}, \opt{c}, \opt{b}) and lowered by $\delta$, keeping the outer
  dimensions (height and depth) of the box. Thus the float placement of \LaTeX{}
  is not affected, and the mark of double-column floats (depth 1\,sp) is kept.
\item[Top floats (\cs{@cflt})] $p=\opt{t}$, $\delta=\cs{topskip}-h$; the top of the
  area is on the top of the character frame of the first line.
\item[Bottom floats (\cs{@cflb})] $p=\opt{b}$, $\delta=d-d_{\mathrm{last}}$, where
  $d_{\mathrm{last}}$ is the depth of the last box in the column modulo $b$ (it is
  $d$ if the last box is a gyou-dori area). After the floats, an empty line of
  height 0 and depth $d$ and glue of ``$r-d$'' are added, where $r$ is ``the
  column height $-\cs{topskip}$'' modulo $b$. Then, even when the floats are
  pushed down by the \cs{vfil} of \cs{clearpage}, their bottom is on the bottom of
  the character frame of the last line in the column, and no overfull box occurs
  in classes where that bottom goes below the type area. When pushed down (the
  column ends with glue), $\delta=0$ and the stretch of \cs{@textbottom} is
  cancelled.
\item[Here floats (\cs{@addtocurcol})] $p=\opt{c}$,
  $\delta=g/2+d-d_{\mathrm{prev}}$, where $d_{\mathrm{prev}}$ is the depth of the
  previous line modulo $b$. A \cs{nobreak} is inserted before the float so that
  it does not start a column.
\item[Rounding of separations] Each separation is set to the smallest value of
  the form $g+kb$ ($k\ge0$ an integer) that is not less than the given value,
  without stretch or shrink; \cs{intextsep}, used above and below, is
  $g/2+kb$. Then ``float area $+$ separation'' is a multiple of $b$, and the text
  lines are not shifted.
\item[Order of footnotes] The \cs{@makecol} of pLaTeX puts footnotes after the
  floats, so their height is subtracted from the column height when computing
  $r$. The \cs{@makecol} of \LaTeX{} 2025-06-01 or later (e.g.\ LuaLaTeX) puts the
  floats after the footnotes, so in a column with bottom floats a \cs{vfil} is put
  before the footnotes to place the footnotes and the floats at the bottom.
\end{description}

\subsection{Interfaces to other packages}

\begin{description}
\item[Double-column floats] They are handled in \cs{@comdblflelt}. A float in
  \cs{@dbltoplist} is at the top: $p=\opt{t}$, $\delta=\cs{topskip}-h$. Otherwise it is at the bottom
  (\pkg{nidanfloat}, \pkg{stfloats}): $p=\opt{b}$, $\delta=d-r$ with $r$ being
  ``$\cs{textheight}-\cs{topskip}$'' modulo $b$, and a \cs{vfil} is put before the
  floats.
\item[\pkg{nidanfloat}] Top floats of each column are handled as by \cs{@cflt}.
  \cs{L@cflb} and \cs{R@cflb} are replaced so that the bottom floats are aligned to
  the last line in the column, with the stretch between the text and the floats.
\item[\pkg{flushend}] After the columns of the last page are built, they are
  rebuilt without stretching or shrinking the glue in them (with fill glue at the
  end). Where the two columns are joined (\cs{atColsBreak}), the distance from
  the baseline of the last element of the left column to that of the first
  element of the right column is rounded up to a multiple of $b$, taking the
  shift $\delta$ of float boxes into account.
\item[\pkg{evenend}] Before the copies of column floats on the last page are
  taken, top floats are realigned to \opt{t} and bottom floats to \opt{b}.
\end{description}

\section{Limitations}

\begin{itemize}
\item Gyou-dori of display math and floats is based on the line pitch of
  \cs{normalsize}; inside \cs{small} etc.\ the lines do not match the
  surrounding text. The \env{gyoudori} environment and \cs{gyoudoriarea} keep
  the surrounding lines on their grid with \opt{local}
  (Section~\ref{sec:env}).
\item \cs{topsep}, \cs{itemsep} and \cs{parsep} of lists are not multiples of
  the line pitch, so lines in and around lists are off the grid.
\item Float-only pages and columns are not handled.
\item Vertical writing is not supported. \cs{flushbottom} and the \opt{[H]}
  placement of the \pkg{float} package have not been tested.
\item When \pkg{flushend}, footnotes, and a bottom or here float on the last page
  come together, text may be shifted by about 1\,pt.
\item \cs{fnbelowfloat} of \pkg{stfloats} has no effect in pLaTeX, which uses
  its own \cs{@makecol}.
\item In classes whose text height is not on the line grid (e.g.\ \pkg{jarticle}
  of pLaTeX), there is a remainder below the last line, so the area of display
  math may extend below the last line by up to one line. The \env{gyoudori}
  environment and \cs{gyoudoriarea} allow no break right after them, so they
  move to the next page (column) instead.
\item Known, not yet fixed: with \pkg{nidanfloat} and \opt{displaymath} together,
  text lines may be shifted off the grid. The processing of \pkg{nidanfloat} that
  moves text overflowing the left column into the right column seems to interfere
  with the large depth of the boxes of gyou-dori areas.
\end{itemize}

\section{References}

\begin{itemize}
\item W3C, \textit{Requirements for Japanese Text Layout},
  \url{https://www.w3.org/TR/jlreq/}
\item Noriyuki Abe, \textit{Gyodori} (in Japanese),
  \url{https://abenori.blogspot.com/2017/12/gyodori.html}
\end{itemize}

\section{License}

MIT License.

\end{document}
